\uuid{27ge}
\exo7id{4588}
\titre{exo7 4588}
\auteur{quercia}
\organisation{exo7}
\datecreate{2010-03-14}
\isIndication{false}
\isCorrection{true}
\chapitre{Série entière}
\sousChapitre{Calcul de la somme d'une série entière}
\cpge{mpsi-mp2i.nombres-reels-suites.d, pcsi.nombres-reels-suites.c, ptsi.nombres-reels-suites.c, mpsi-mp2i.analyse-asymptotique.b, pcsi.analyse-asymptotique.b, ptsi.analyse-asymptotique.b, bcpst1.suites-reelles, mp-mpi.series-entieres.a, pc.series-entieres.a, psi.series-entieres.a, pt.series-entieres.a, mp-mpi.series-entieres.d, pc.series-entieres.c, psi.series-entieres.c, pt.series-entieres.c}
\module{Analyse}
\niveau{L2}
\difficulte{}

\contenu{
\question{
Soit $(u_n)$ définie par, pour tout $n\in \N$, $\sum_{k=0}^n\frac{u_{n-k}}{k!\strut}=1$. Trouver la
limite de $(u_n)$.
}
\reponse{
$f(t) = \sum_{n=0}^\infty u_nt^n = \frac{e^{-t}}{1-t\strut} = \sum_{n=0}^\infty\sum_{k=0}^n\frac{(-1)^k}{\strut k!}t^n$
donc $u_n\to\frac1{\strut e}$ lorsque $n\to\infty$.
}
}
